% Part: propositional-logic
% Chapter: syntax-and-semantics
% Section: formulas

\documentclass[../../../include/open-logic-section]{subfiles}

\begin{document}

\olfileid{pl}{syn}{fml}

\olsection{ప్రతిజ్ఞావాక్య \usetoken{P}{formula}}

ప్రతిజ్ఞావాక్యాత్మక తర్కంలోని \tetoken{సూత్రాలను}{!!^{formula}s} \emph{\tetoken{ప్రతిజ్ఞావాక్య చరాలు}{!!{propositional
    variable}s}}\iftag{prvFalse}{\iftag{prvTrue}{,}{ మరియు} ప్రతిజ్ఞావాక్య
  స్థిరాంకం~$\lfalse$}{}\iftag{prvTrue}{\iftag{prvFalse}{
    మరియు}{} ప్రతిజ్ఞావాక్య స్థిరాంకం~$\ltrue$}{} నుంచి \emph{తార్కిక
  సంయోజకాలను} ఉపయోగించి నిర్మిస్తాం.

\begin{enumerate}
\item \tetoken{ప్రతిజ్ఞావాక్య చరాలు}{!!{propositional variable}s} $\Obj p_0$, $\Obj p_1$, \dots గల
  \tetoken{ఒక అనంతంగా లెక్కించదగిన}{!!^a{denumerable}} సమితి~$\PVar$
\tagitem{prvFalse}{\tetoken{అసత్యం}{!!{falsity}} కోసం ప్రతిజ్ఞావాక్య స్థిరాంకం~$\lfalse$.}{}
\tagitem{prvTrue}{\tetoken{సత్యం}{!!{truth}} కోసం ప్రతిజ్ఞావాక్య స్థిరాంకం~$\ltrue$.}{}
\item తార్కిక సంయోజకాలు:
  \startycommalist
  \iftag{prvNot}{\ycomma $\lnot$ (నిషేధం)}{}%
  \iftag{prvAnd}{\ycomma $\land$ (సంయోగం)}{}%
  \iftag{prvOr}{\ycomma $\lor$ (వికల్పం)}{}%
  \iftag{prvIf}{\ycomma $\lif$ (\tetoken{సోపాధికం}{!!{conditional}})}{}%
  \iftag{prvIff}{\ycomma $\liff$ (\tetoken{ద్విసోపాధికం}{!!{biconditional}})}{}%
\item విరామచిహ్నాలు: (, ), మరియు కామా.
\end{enumerate}

ప్రతిజ్ఞావాక్యాత్మక తర్కంలోని ఈ భాషను $\Lang L_0$తో సూచిస్తాం.

\iftag{defNot,defOr,defAnd,defIf,defIff,defTrue,defFalse,defEx,defAll}{%
పైన పరిచయం చేసిన ప్రాథమిక సంయోజకాలతో పాటు కింది
\emph{నిర్వచిత} సంకేతాలను కూడా వాడతాం:
\startycommalist
  \iftag{defNot}{\ycomma $\lnot$ (నిషేధం)}{}%
  \iftag{defAnd}{\ycomma $\land$ (సంయోగం)}{}%
  \iftag{defOr}{\ycomma $\lor$ (వికల్పం)}{}%
  \iftag{defIf}{\ycomma $\lif$ (\tetoken{సోపాధికం}{!!{conditional}})}{}%
  \iftag{defIff}{\ycomma $\liff$ (\tetoken{ద్విసోపాధికం}{!!{biconditional}})}{}%
  \iftag{defFalse}{\ycomma $\lfalse$ (\tetoken{అసత్యం}{!!{falsity}})}{}%
  \iftag{defTrue}{\ycomma $\ltrue$ (\tetoken{సత్యం}{!!{truth}})}{}.
}{}
\sourcecorrection{OLTEPLSYN-001}{ఆంగ్ల మూలంలోని చివరి defTrue ట్యాగు
పరీక్షకు ఖాళీ అసత్య శాఖ రాకముందే బయటి ట్యాగు సమూహపు ముగింపు కుండలీకరణ
వచ్చింది. లోపలి పరీక్షను రెండు శాఖలతో ముగించి, తరువాత బయటి పరీక్షను
ముగిసేలా కుండలీకరణల క్రమాన్ని సరిచేశాం; పాఠ్య విషయం మారలేదు.}

\begin{tagblock}{defNot,defOr,defAnd,defIf,defIff,defTrue,defFalse,defEx,defAll}
\begin{explain}
నిర్వచిత సంకేతం అధికారికంగా భాషలో భాగం కాదు; దాన్ని అనౌపచారిక
సంక్షిప్తీకరణగా పరిచయం చేస్తాం. ప్రాథమిక సంకేతాలను మాత్రమే వాడితే
చాలా పొడవయ్యే సూత్రాలను సంక్షిప్తంగా రాయడానికి అది వీలు కల్పిస్తుంది.
ఇది స్పష్టమైన ప్రయోజనం. అయితే ఇంకా పెద్ద ప్రయోజనం ఏమిటంటే నిరూపణలు
చిన్నవవుతాయి. ఒక సంకేతం ప్రాథమికమైతే నిరూపణల్లో దాన్ని విడిగా
నిర్వహించాలి. అందువల్ల ప్రాథమిక సంకేతాలు ఎక్కువైతే నిరూపణలు
పొడవవుతాయి.
\end{explain}
\end{tagblock}

% Alternate symbols

\begin{intro}
మేము పైన వాడిన వాటికి భిన్నమైన పరిభాష, సంకేతాలు మీకు పరిచయం
ఉండవచ్చు. తర్కశాస్త్ర గ్రంథాలు (మరియు ఉపాధ్యాయులు) “నిషేధం” కోసం
సాధారణంగా $\sim$, $\neg$, మరియు~!\ను; “సంయోగం” కోసం $\wedge$,
$\cdot$, మరియు $\&$ను వాడుతారు. “షరతీయ సంయోజకం” లేదా “అన్వయం” కోసం
సాధారణంగా వాడే సంకేతాలు $\rightarrow$, $\Rightarrow$, మరియు
$\supset$.
\iftag{prvIff,defIff}{“ద్విషరతీయ సంయోజకం,” “ద్వ్యన్వయం,” లేదా
  “(భౌతిక) తుల్యత” కోసం $\leftrightarrow$, $\Leftrightarrow$,
  మరియు $\equiv$ సంకేతాలు వాడతారు.}{}
\iftag{prvFalse,defFalse}{$\lfalse$ సంకేతాన్ని “అసత్యం,” “ఫాల్సమ్,”
  “అసంభవం,” లేదా “బాటమ్” అని వేర్వేరుగా పిలుస్తారు.}{}
\iftag{prvTrue,defTrue}{$\ltrue$ సంకేతాన్ని “సత్యం,” “వేరమ్,” లేదా
  “టాప్” అని వేర్వేరుగా పిలుస్తారు.}{}
\end{intro}

\begin{defn}[సూత్రం]
\ollabel{defn:formulas}
ప్రతిజ్ఞావాక్యాత్మక తర్కంలోని \emph{\tetoken{సూత్రాలు}{!!{formula}s}} సమితి~$\Frm[L_0]$ను కింది
విధంగా ఆగమనాత్మకంగా నిర్వచిస్తాం:
\begin{enumerate}
\tagitem{prvFalse}{$\lfalse$ ఒక పరమాణు \tetoken{సూత్రం}{!!{formula}}.}{}

\tagitem{prvTrue}{$\ltrue$ ఒక పరమాణు \tetoken{సూత్రం}{!!{formula}}.}{}

\item ప్రతి \tetoken{ప్రతిజ్ఞావాక్య చరం}{!!{propositional variable}}~$\Obj p_i$ ఒక పరమాణు
  \tetoken{సూత్రం}{!!{formula}}.

\tagitem{prvNot}{$!A$ ఒక \tetoken{సూత్రం}{!!a{formula}} అయితే, $\lnot !A$ ఒక
  \tetoken{సూత్రం}{!!a{formula}}.}{}

\tagitem{prvAnd}{$!A$, $!B$లు \tetoken{సూత్రాలు}{!!{formula}s} అయితే, $(!A \land
  !B)$ ఒక \tetoken{సూత్రం}{!!a{formula}}.}{}

\tagitem{prvOr}{$!A$, $!B$లు \tetoken{సూత్రాలు}{!!{formula}s} అయితే, $(!A \lor !B)$
  ఒక \tetoken{సూత్రం}{!!a{formula}}.}{}

\tagitem{prvIf}{$!A$, $!B$లు \tetoken{సూత్రాలు}{!!{formula}s} అయితే, $(!A \lif !B)$
  ఒక \tetoken{సూత్రం}{!!a{formula}}.}{}

\tagitem{prvIff}{$!A$, $!B$లు \tetoken{సూత్రాలు}{!!{formula}s} అయితే, $(!A \liff !B)$
  ఒక \tetoken{సూత్రం}{!!a{formula}}.}{}

\tagitem{limitClause}{మరేదీ \tetoken{ఒక సూత్రం}{!!a{formula}} కాదు.}{}
\end{enumerate}
\end{defn}

\begin{explain}
\tetoken{సూత్రాలు}{!!{formula}s} నిర్వచనం ఒక \emph{ఆగమనాత్మక నిర్వచనం}. సారాంశంగా,
\tetoken{సూత్రాలు}{!!{formula}s} సమితిని అనంత దశల్లో నిర్మిస్తాం. మొదటి దశలో అన్ని
పరమాణు సూత్రాలూ సూత్రాలని ప్రకటిస్తాం; ఇది నిర్వచనంలోని మొదటి కొన్ని
సందర్భాలకు, అంటే
\iftag{prvTrue}{$\ltrue$, }{}%
\iftag{prvFalse}{$\lfalse$, }{}%
$\Obj p_i$ సందర్భాలకు సరిపోతుంది. కాబట్టి “పరమాణు \tetoken{సూత్రం}{!!{formula}}”
అంటే ఈ రూపాల్లో ఏదో ఒక రూపంలో ఉన్న \tetoken{సూత్రం}{!!{formula}}.

నిర్వచనంలోని మిగతా సందర్భాలు, అప్పటికే నిర్మించిన \tetoken{సూత్రాలు}{!!{formula}s}
నుంచి కొత్త \tetoken{సూత్రాలను}{!!{formula}s} నిర్మించే నియమాలను ఇస్తాయి. రెండవ దశలో
వాటిని ఉపయోగించి పరమాణు \tetoken{సూత్రాలు}{!!{formula}s} నుంచి \tetoken{సూత్రాలను}{!!{formula}s}
నిర్మించవచ్చు. మూడవ దశలో పరమాణు సూత్రాలు, రెండవ దశలో పొందిన
సూత్రాల నుంచి కొత్త సూత్రాలను నిర్మిస్తాం; ఇలాగే కొనసాగిస్తాం.
అటువంటి ఏదో ఒక దశలో చివరికి నిర్మితమయ్యేదే \tetoken{సూత్రం}{!!{formula}}; మరేదీ
సూత్రం కాదు.
\end{explain}

$!B$, $!C$ల నుంచి ద్విస్థాన సంయోజకం~$\ast$ను ఉపయోగించి నిర్మించిన
$(!B \ast !C)$ అనే సూత్రాన్ని రాసేటప్పుడు, బయటివైపు ఉన్న
కుండలీకరణల జతను తరచుగా వదిలి, కేవలం~$!B \ast !C$ అని రాస్తాం.

\begin{tagblock}{defNot,defOr,defAnd,defIf,defIff,defTrue,defFalse,defEx,defAll}
\begin{defn}
నిర్వచిత సంయోజకాలతో నిర్మించిన సూత్రాలను కింది విధంగా అర్థం
చేసుకోవాలి:

\begin{tagenumerate}{defTrue,defFalse,defNot,defOr,defAnd,defIf,defIff,defEx,defAll}

\tagitem{defTrue}{$\ltrue$ అనేది
  \iftag{prvFalse}{$\lnot\lfalse$}{ఏదో స్థిర పరమాణు
    \tetoken{సూత్రం}{!!{formula}}~$!A$కు $(!A \lor \lnot !A)$}కు సంక్షిప్త రూపం.}{}

\tagitem{defFalse}{$\lfalse$ అనేది
  \iftag{prvTrue}{$\lnot\ltrue$}{ఏదో స్థిర పరమాణు
    \tetoken{సూత్రం}{!!{formula}}~$!A$కు $(!A \land \lnot !A)$}కు సంక్షిప్త రూపం.}{}

\tagitem{defNot}{$\lnot !A$ అనేది $!A \lif \lfalse$కు సంక్షిప్త
  రూపం.}{}

\tagitem{defOr}{$!A \lor !B$ అనేది
  \iftag{prvAnd}{$\lnot(\lnot !A \land \lnot !B)$}{$\lnot !A \lif
    !B$}కు సంక్షిప్త రూపం.}{}

\tagitem{defAnd}{$!A \land !B$ అనేది
  \iftag{prvOr}{$\lnot(\lnot !A \lor \lnot !B)$}{$\lnot (!A \lif
    \lnot !B)$}కు సంక్షిప్త రూపం.}{}

\tagitem{defIf}{$!A \lif !B$ అనేది
  \iftag{prvOr}{$\lnot !A \lor !B$}{$\lnot (!A \land \lnot !B)$}కు
  సంక్షిప్త రూపం.}{}
\sourcecorrection{OLTEPLSYN-002}{ఆంగ్ల మూలంలోని మొదటి సంక్షిప్త
సూత్రం చివర సరిపోలే ఎడమ కుండలీకరణ లేని కుడి కుండలీకరణ ఉంది. బయటివైపు
కుండలీకరణలను వదలవచ్చనే ముందు నియమానికి అనుగుణంగా ఆ అదనపు కుడి
కుండలీకరణను మాత్రమే తొలగించాం.}

\tagitem{defIff}{$!A \liff !B$ అనేది $(!A \lif !B) \land (!B
  \lif !A)$కు సంక్షిప్త రూపం.}{}
\end{tagenumerate}
\end{defn}
\end{tagblock}

\begin{defn}[వాక్యనిర్మాణ తాదాత్మ్యం]
$\ident$ సంకేతం సంకేతమాలల మధ్య వాక్యనిర్మాణ తాదాత్మ్యాన్ని
వ్యక్తపరుస్తుంది; అంటే, $!A \ident !B$ కావడం, $!A$, $!B$లు ఒకే
పొడవు గల సంకేతమాలలై, ప్రతి స్థానంలో ఒకే సంకేతాన్ని కలిగి
ఉన్నప్పుడూ, అప్పుడు మాత్రమే.
\end{defn}

$\ident$ సంకేతానికి ఇరువైపులా సంధానం ద్వారా పొందిన సంకేతమాలలు
ఉండవచ్చు. ఉదాహరణకు, $!A \ident (!B \lor !C)$ అంటే: తెరచే
కుండలీకరణ, $!B$ అనే సంకేతమాల, $\lor$ సంకేతం, $!C$ అనే
సంకేతమాల, మూసే కుండలీకరణలను ఈ క్రమంలో సంధానించి పొందిన సంకేతమాలే
$!A$. అలా అయితే $!A$లోని మొదటి సంకేతం తెరచే కుండలీకరణ అని, రెండవ
సంకేతం వద్ద మొదలయ్యే ఉపమాలగా $!B$ను $!A$ కలిగి ఉందని, ఆ ఉపమాల
తరువాత $\lor$ వస్తుందని, తదితర విషయాలను తెలుసుకుంటాం.

\end{document}
